% !PRINTED 782
\weberpage{782}
On the other hand, however, by 2.:

\weq{19}{(s-1) Q_1 Q_2 \ldots Q_\mu}

is finite for $s = 1$, and so if we divide (19) by (18), it follows that:
\[
(s-1)^{1-\frac{\mu}{\nu}} \text{ finite for } s = 1 \text{,}
\]
i.e.

\weq{20}{\nu \geqq \mu.}

Taken together with (14), however, this yields $\nu = \mu$, whereby the following theorem is proved:

3. \gesperrt{The degree of the cyclotomic field $\Omega_m$ is equal to $\varphi(m)$, and hence the cyclotomic equation of degree $\varphi(m)$, whose roots are the primitive $m$th roots of unity, is irreducible.}

But this consideration renders us yet another important service, in that it furnishes us with the proof of the famous theorem \gesperrt{that in every arithmetic progression whose difference and initial term are natural numbers without a common divisor there are contained infinitely many primes}\textsuperscript{1)}.

For it now follows from (18) that the product
\[
(s-1) Q_1 Q_2 \ldots Q_\mu
\]
has a finite value different from zero for $s = 1$, and since by 2. none of the factors is infinite, it follows that:

4. \gesperrt{The sums}
\[
(s-1) Q_1(s), Q_2(s), \ldots, Q_\mu(s)
\]
\gesperrt{have finite values different from zero for $s = 1$.}

If we apply to all the factors of (11) the formula:
\[
\log \left[1 - \chi(p) p^{-s}\right] = \frac{\chi(p)}{p^s} + \frac{1}{2} \frac{\chi(p^2)}{p^{2s}} + \frac{1}{3} \frac{\chi(p^3)}{p^{3s}} + \cdots % [sic] the source omits the mathematically required minus sign before \log
\]

\noindent\rule{4cm}{0.4pt}

\textsuperscript{1)} The first complete and general proof of this theorem was given by \gesperrt{Dirichlet} (Abhandlungen der Berl. Akademie vom Jahre 1837. Gesammelte Werke Nr.\ XXI). \gesperrt{Dedekind} has drawn attention to the essential simplification of this proof by the use of the \gesperrt{Kummer} formulae for the class number in the cyclotomic fields (Vorlesungen über Zahlentheorie, 3. Auflage, S. 596).
%CONT>>
